\documentclass[11pt]{article}
\usepackage[margin=1in]{geometry}
\usepackage{amsmath,amssymb,amsthm}
\usepackage{graphicx}
\usepackage{hyperref}

\newcommand{\M}{\mathcal M}
\newcommand{\h}{\mathsf h}
\newcommand{\bmo}{\mathsf{bmo}}

\newtheorem{theorem}{Theorem}

\title{A one-sided endpoint complex interpolation result for\\
conditioned noncommutative martingale Hardy and BMO spaces}
\author{Automatic Math Discovery Run \texttt{fa\_banach\_001}}
\date{June 15, 2026}

\begin{document}
\maketitle

\section*{Status}

This is a candidate partial result for the open problem in
Randrianantoanina's paper \cite{Ran2021}.  It proves one inclusion in the
range where the target exponent is \(>1\).  It does not prove the reverse
inclusion, and it does not settle the range where the target exponent is
\(\leq 1\).

\section*{Source problem}

The source paper proves that, for finite indices \(0<p,q<\infty\),
\[
[\h_p^c(\M),\h_q^c(\M)]_\theta=\h_r^c(\M),
\qquad
  \frac1r=\frac{1-\theta}{p}+\frac{\theta}{q}.
\]
Immediately after that theorem, the paper asks whether the corresponding
statement remains valid when \(0<p<1\) and the endpoint couple
\((\h_p^c(\M),\bmo^c(\M))\) is used.

\begin{center}
\includegraphics[width=\textwidth]{figures/open_problem_crop.png}
\end{center}

\noindent
Thus the desired full endpoint identity would be
\[
[\h_p^c(\M),\bmo^c(\M)]_\theta=\h_r^c(\M),
\qquad
\frac1r=\frac{1-\theta}{p}.
\]

\section*{Proof intuition}

The source already proves the analogous one-sided inclusion for the larger
capital Hardy/BMO couple by using Wolff's interpolation lemma at the level of
inclusions.  The same interpolation geometry also applies to the conditioned
column spaces.  The finite-index complex theorem gives
\([\h_p^c,\h_r^c]_\eta=\h_1^c\), while the known Banach-range endpoint theorem
of Bekjan--Chen--Perrin--Yin gives
\([\h_1^c,\bmo^c]_\phi=\h_r^c\) when \(r>1\).  Wolff's inclusion lemma then
splices these two identities and yields the inclusion from
\([\h_p^c,\bmo^c]_\theta\) into \(\h_r^c\).  The argument deliberately uses
only the inclusion half of Wolff's mechanism, because full complex
reiteration is not available in this quasi-Banach endpoint setting.

\begin{theorem}
Let \(\M\) be as in \cite{Ran2021}: a semifinite von Neumann algebra with an
increasing filtration admitting trace-preserving conditional expectations.
Let \(0<p<1\) and \(0<\theta<1\), and put
\[
  \frac1r=\frac{1-\theta}{p}.
\]
If \(r>1\), then
\[
[\h_p^c(\M),\bmo^c(\M)]_\theta
  \hookrightarrow \h_r^c(\M)
\]
continuously.
\end{theorem}

\begin{proof}
Choose \(\eta\in(0,1)\) so that
\[
  1=\frac{1-\eta}{p}+\frac{\eta}{r}.
\]
This is possible because \(p<1<r\).  By the finite-index complex
interpolation theorem of \cite[Theorem 3.16]{Ran2021}, applied to the two
finite exponents \(p\) and \(r\), we have
\[
[\h_p^c(\M),\h_r^c(\M)]_\eta=\h_1^c(\M)
\]
with equivalent quasi-norms.

Set
\[
  \phi=1-\frac1r.
\]
Since \(r>1\), the Banach-range endpoint theorem of
Bekjan--Chen--Perrin--Yin \cite[Theorem 4.4]{BCPY2009}, with the order of
the couple reversed, gives
\[
[\h_1^c(\M),\bmo^c(\M)]_\phi=\h_r^c(\M).
\]

We also need the intersection condition in Wolff's inclusion lemma.  If
\(x\in \h_p^c(\M)\cap\bmo^c(\M)\), then the real endpoint identity
\[
(\h_p^c(\M),\bmo^c(\M))_{\alpha,s}=\h_s^c(\M),
\qquad
\frac1s=\frac{1-\alpha}{p},
\]
from \cite[Theorem 3.5]{Ran2021} implies, by the elementary
\(K\)-functional estimate for elements of the intersection, that
\[
x\in \h_1^c(\M)\cap\h_r^c(\M).
\]
Hence
\[
\h_p^c(\M)\cap\bmo^c(\M)\subset \h_1^c(\M)\cap\h_r^c(\M).
\]

Apply Wolff's complex interpolation lemma in the inclusion form used in
\cite[Corollary 3.19]{Ran2021} to
\[
B_1=\h_p^c(\M),\qquad
B_2=\h_1^c(\M),\qquad
B_3=\h_r^c(\M),\qquad
B_4=\bmo^c(\M).
\]
The two interpolation identities above give the hypotheses
\([B_1,B_3]_\eta\subset B_2\) and \([B_2,B_4]_\phi\subset B_3\).  Therefore
\[
[B_1,B_4]_\Theta\subset B_3,
\qquad
\Theta=\frac{\phi}{1-\eta+\eta\phi}.
\]
It remains only to check the parameter.  Write \(a=1/p\) and \(b=1/r\).  Then
\[
\eta=\frac{a-1}{a-b},\qquad \phi=1-b,
\]
and hence
\[
\Theta
=\frac{1-b}{1-\eta b}
=\frac{1-b}{1-\frac{b(a-1)}{a-b}}
=\frac{a-b}{a}
=1-\frac{p}{r}
=\theta,
\]
because \(1/r=(1-\theta)/p\).  Thus
\[
[\h_p^c(\M),\bmo^c(\M)]_\theta\subset \h_r^c(\M),
\]
as claimed.
\end{proof}

\section*{Limitations}

The proof gives only the inclusion
\[
[\h_p^c(\M),\bmo^c(\M)]_\theta\subset \h_r^c(\M)
\]
when \(r>1\).  It does not give
\[
\h_r^c(\M)\subset [\h_p^c(\M),\bmo^c(\M)]_\theta.
\]
Obtaining that reverse inclusion would require a new endpoint construction
or a full quasi-Banach complex reiteration principle.  This is precisely the
obstruction identified in the source paper.  The argument also does not cover
the case \(r\leq1\), since the required intermediate Banach-range theorem
\([\h_1^c,\bmo^c]_\phi=\h_r^c\) is only available for \(r>1\).

\section*{Verification notes}

The proof is symbolic and uses no numerical computation.  The parameter
calculation was checked algebraically in the proof.  The main review point is
the use of Wolff's interpolation lemma in the same quasi-Banach inclusion
form invoked by \cite[Corollary 3.19]{Ran2021}.

\section*{Novelty check}

The run's lightweight indexes were searched for arXiv:2108.06341, title
terms, and core phrases such as \(\h_p^c\), \(\bmo^c\), and complex
interpolation; no prior packet or attempt was found.  Web searches for exact
and close variants of the endpoint question found the source paper, the
Bekjan--Chen--Perrin--Yin Banach-range theorem \cite{BCPY2009}, the 2022
real-method \( \h_\infty^c \) sequel, and the 2024 capital-H real-method
sequel, but no later full solution of this conditioned BMO complex endpoint.

\begin{thebibliography}{9}

\bibitem{Ran2021}
N.~Randrianantoanina,
\emph{Interpolation between noncommutative martingale Hardy and BMO spaces:
the case \(0<p<1\)},
arXiv:2108.06341.

\bibitem{BCPY2009}
T.~N. Bekjan, Z.~Chen, M.~Perrin, and Z.~Yin,
\emph{Atomic decomposition and interpolation for Hardy spaces of
noncommutative martingales},
arXiv:0906.4437.

\bibitem{Wolff}
T.~Wolff,
\emph{A note on interpolation spaces},
Harmonic Analysis (Minneapolis, Minn., 1981), Lecture Notes in Math. 908,
Springer, 1982, 199--204.

\end{thebibliography}

\end{document}
